%_-------------------------
% Resume in Latex
% Author : Jake Gutierrez
% Based off of: https://github.com/sb2nov/resume
% License : MIT
%------------------------

\documentclass[letterpaper,11pt]{article}

\usepackage{latexsym}
\usepackage[empty]{fullpage}
\usepackage{titlesec}
\usepackage{marvosym}
\usepackage[usenames,dvipsnames]{color}
\usepackage{verbatim}
\usepackage{enumitem}
\usepackage[hidelinks]{hyperref}
\usepackage{fancyhdr}
\usepackage[english, russian]{babel}
\usepackage{tabularx}
\input{glyphtounicode}

\usepackage[T2A]{fontenc}
\usepackage[utf8]{inputenc}
\usepackage[russian]{babel}
%----------FONT OPTIONS----------
% sans-serif
% \usepackage[sfdefault]{FiraSans}
% \usepackage[sfdefault]{roboto}
% \usepackage[sfdefault]{noto-sans}
% \usepackage[default]{sourcesanspro}

% serif
% \usepackage{CormorantGaramond}
% \usepackage{charter}


\pagestyle{fancy}
\fancyhf{} % clear all header and footer fields
\fancyfoot{}
\renewcommand{\headrulewidth}{0pt}
\renewcommand{\footrulewidth}{0pt}

% Adjust margins
\addtolength{\oddsidemargin}{-0.5in}
\addtolength{\evensidemargin}{-0.5in}
\addtolength{\textwidth}{1in}
\addtolength{\topmargin}{-.5in}
\addtolength{\textheight}{1.0in}

\urlstyle{same}

\raggedbottom
\raggedright
\setlength{\tabcolsep}{0in}

% Sections formatting
\titleformat{\section}{
  \vspace{-4pt}\scshape\raggedright\large
}{}{0em}{}[\color{black}\titlerule \vspace{-5pt}]

% Ensure that generate pdf is machine readable/ATS parsable
\pdfgentounicode=1


%-----------CUSTOM-COMMANDS-----------
\newcommand{\resumeItem}[1]{
  \item\small{
    {#1 \vspace{-2pt}}
  }
}

\newcommand{\resumeSubheading}[4]{
  \vspace{-2pt}\item
    \begin{tabular*}{0.97\textwidth}[t]{l@{\extracolsep{\fill}}r}
      \textbf{#1} & #2 \\
      \textit{\small#3} & \textit{\small #4} \\
    \end{tabular*}\vspace{-7pt}
}

\newcommand{\resumeSubheadingExperience}[2]{
  \vspace{-2pt}\item
    \begin{tabular*}{0.97\textwidth}[t]{l@{\extracolsep{\fill}}r}
      \textbf{#1} & #2 \\
    \end{tabular*}\vspace{-7pt}
}

\newcommand{\resumeSubSubheading}[2]{
    \item
    \begin{tabular*}{0.97\textwidth}{l@{\extracolsep{\fill}}r}
      \textit{\small#1} & \textit{\small #2} \\
    \end{tabular*}\vspace{-7pt}
}

\newcommand{\resumeProjectHeading}[2]{
    \item
    \begin{tabular*}{0.97\textwidth}{l@{\extracolsep{\fill}}r}
      \small#1 & #2 \\
    \end{tabular*}\vspace{-7pt}
}

\newcommand{\resumeSubItem}[1]{\resumeItem{#1}\vspace{-4pt}}

\renewcommand\labelitemii{$\vcenter{\hbox{\tiny$\bullet$}}$}

\newcommand{\resumeSubHeadingListStart}{\begin{itemize}[leftmargin=0.15in, label={}]}
\newcommand{\resumeSubHeadingListEnd}{\end{itemize}}
\newcommand{\resumeItemListStart}{\begin{itemize}}
\newcommand{\resumeItemListEnd}{\end{itemize}\vspace{-5pt}}

%-------------------------------------------
%%%%%%  RESUME STARTS HERE  %%%%%%%%%%%%%%%%%%%%%%%%%%%%
\begin{document}


%----------ЗАГОЛОВОК----------
\begin{center}
    \textbf{\Huge \scshape Фамилия Имя} \\ \vspace{1pt}
    \small телефон $|$ \href{mailto:x@x.com}{\underline{почта}} $|$ 
    \href{https://t.me/}{\underline{телеграмм через @}} $|$
    \href{https://github.com/}{\underline{github.com/<your nickname>}} \\ \vspace{2pt}
    \largeСпециалист по Информационной Безопасности
\end{center}


%----------О-СЕБЕ----------
\section{О Себе}
 \begin{itemize}[leftmargin=0.15in, label={}]
    \small{\item{
     Студент с 4+ годами опыта программирования на Python и 2-летним опытом работы с Arch Linux и C++. Участвовал в CTF-соревнованиях, финалист олимпиады Газпром, МОШ (аэрокосмос) и призёр Национальной Технологической Олимпиады по Информационной Безопасности. Разрабатывал проекты, включая телеграм-ботов и инструмент для модификации оперативной памяти приложений с использованием DLL-инъекций. Интересуюсь низкоуровневым программированием, работой процессоров и операционных систем. Стремлюсь развиваться в сфере информационной безопасности, применяя свои знания в программировании и математике.  
    }} 
 \end{itemize}


%-----------ОБРАЗОВАНИЕ-----------
\section{Образование}
  \resumeSubHeadingListStart
    \resumeSubheading
      {Центральный Университет}{Москва}
      {Математика и компьютерные науки, трек Разработка}{2024 - 2028}
    \resumeItemListStart
      \resumeItem{Релевантные курсы: C++, Линейная алгебра, Дискретная математика, Основы промышленного программирования (Java, Go)}
      \resumeItem{Стипендия, покрывающая 50\% стоимости обучения}
    \resumeItemListEnd
  \resumeSubHeadingListEnd


%-----------ПРОЕКТЫ-----------
\section{Проекты}
    \resumeSubHeadingListStart
      \resumeProjectHeading
          {\textbf{Чит для игры Control} $|$ \emph{C++, анализ и модификация оперативной памяти, DLL-инъекция}}{Апрель 2023}
          \resumeItemListStart
            \resumeItem{Провел анализ памяти игры и выявил основные оффсеты}
            \resumeItem{Использовал Cheat Engine - отладчик и инструмент для reverse engineering}
            \resumeItem{Использовал DLL-инъекцию для интеграции программы в процесс игры}
            \resumeItem{Углубил понимание работы оперативной памяти, низкоуровнего программирования и процессов в Windows}
          \resumeItemListEnd
      \resumeProjectHeading
          {\textbf{Телеграм-бот для рекомендаций фильмов} $|$ \emph{aiogram3, KinopoiskAPI, SQLAlchemy}}{Июль 2023 - сейчас}
          \resumeItemListStart
            \resumeItem{Разработал телеграм-бота для рекомендаций фильмов для антикинотеатра Kinoroom}
            \resumeItem{Использовал KinopoiskAPI для получения актуальной информации по фильмам}
            \resumeItem{Реализовал фильтрации и категоризацию фильмов по разным параметрам}
            \resumeItem{Использовал SQLAlchemy вместе с SQLite, а позже с PostgreSQL}
          \resumeItemListEnd
      \resumeProjectHeading
          {\textbf{Судоку} $|$ \emph{Python, PyQt5, Git, алгоритмы}}{Ноябрь 2023}
          \resumeItemListStart
            \resumeItem{Разработал игру Судоку с использованием PyQt5 для создания графического интерфейса.}
            \resumeItem{Реализовал алгоритм генерации уникальных игровых полей с различными уровнями сложности}
            \resumeItem{Добавил функционал для очистки поля, подсказок и системы подсчёта опыта}
            \resumeItem{Использовал Git для структурированной и безопасной разработки}
          \resumeItemListEnd
    \resumeSubHeadingListEnd


%-----------ОПЫТ-----------
\section{Достижения}
  \resumeSubHeadingListStart
    \resumeSubheadingExperience
      {Национальная Технологическая Олимпиада по Информационной Безопасности}{2024}
      \resumeItemListStart
        \resumeItem{Расследовал инцидент произошедший в компании, исследовал Windows образ на предмет следов злоумышленника}
        \resumeItem{Нашел ключ шифрования данных с помощью которого были зашифрованы данные на образе}
        \resumeItem{Написал отчет о решенных Jeopardy-based заданиях и расследовании инцидента}
      \resumeItemListEnd
    
    \resumeSubheadingExperience
      {Московская Олимпиада Школьников по аэрокосмосу}{2024}
    \resumeItemListStart
      \resumeItem{Написал программу для передачи и шифрования данных для разных модулей и частот}
      \resumeItem{Собрал модель спутника и наземной станции на основе Arduino}
    \resumeItemListEnd
  \resumeSubHeadingListEnd


%-----------НАВЫКИ-----------
\section{Технические навыки}
 \begin{itemize}[leftmargin=0.15in, label={}]
    \small{\item{
     \textbf{Языки программирования}{: Python (продвинутый), C/C++ (средний), SQL (SQLite, PostgreSQL) (средний), JavaScript (начальный), HTML/CSS (начальный), Java (начальный), Go (начальный), bash} \\
     \textbf{Фреймворки}{: Flask, aiogram3} \\
     \textbf{Инструменты}{: Git, Docker, VS Code, PyCharm, CLion} \\
     \textbf{Инструменты}{: nmap, wireshark, volatility, ssh, burpsuite} \\
     \textbf{Библиотеки}{: pandas, NumPy, Matplotlib, requests, BeautifulSoup4, PyQt5} \\ 
     \textbf{Языки}{: Русский (родной), Английский (B2), Японский (начальный)}
    }} \\
 \end{itemize}


%-------------------------------------------
\end{document}